%%=============================================================================
%% Conclusie
%%=============================================================================

\chapter{Conclusie}
\label{ch:conclusie}
Deze bachelorproef onderzocht hoe de kloof tussen de wetenschappelijke logica van het EMAV-project bij ILVO en de technische implementatie in C\#/.NET kon worden verkleind. 
De centrale vaststelling is dat de huidige situatie ertoe leidt dat onderzoekers het systeem als een “black box” ervaren: de berekeningen zijn beschikbaar, maar de onderliggende intentie, aannames, dataflow en validatiestappen zijn niet steeds duidelijk genoeg gedocumenteerd of reproduceerbaar. 
Daardoor ontstaat documentation drift, neemt de kans op misinterpretatie toe en wordt de ontwikkeling van nieuwe experimenten en analyses vertraagd.

Uit het onderzoek blijkt dat een volledige migratie van het bestaande model naar Python niet de meest realistische of efficiënte oplossing is. 
De huidige .NET-architectuur is te complex en te sterk ingebed in de bedrijfslogica om zonder meer te vervangen. 
De meer praktische en duurzame route is een hybride aanpak: de kern van de wetenschappelijke berekeningen en de domeinkennis blijven in de bestaande technisch stabiele omgeving, terwijl onderzoekers toegang krijgen via Python, Jupyter-notebooks en gecontroleerde experimentatieomgevingen. 
Deze keuze houdt de stabiliteit van het systeem intact, maar geeft onderzoekers wel meer autonomie, flexibiliteit en inzicht.

De proof-of-concept bevestigt deze conclusie. 
De uitgevoerde experimenten laten zien dat onderdelen van EMAV succesvol kunnen worden benaderd, getest en geanalyseerd in Python, zonder dat de bestaande C\#-implementatie direct hoeft te worden omgevormd. 
Python.NET bleek daarbij een geschikte keuze omdat het moderne Python-ondersteuning biedt en daarmee het data-science ecosysteem toegankelijk maakt. 
Vanaf dat pythonnet met de dependency injection en andere interne zaken van .NET moet werken, zijn er meestal meer nadelen en frustratie dan de voordelen die python kan opleveren. 
Jupyter-notebooks bieden een laagdrempelige manier om berekeningen te reproduceren en te communiceren

Belangrijk is echter dat deze aanpak geen technologische wondermiddel is. 
De grootste uitdaging is niet alleen het kiezen van de juiste tooling, maar ook het structureren van samenwerking en proces. 
De analyse van documentation drift toont aan dat informatie vaak verloren gaat omdat code, documentatie, analyse-werkflows en communicatiekanalen apart van elkaar opereren. 
Daarom is een Single Source of Truth essentieel: een centrale, duidelijke, versiebeheerbare bron waarin de modeldefinities, aannames, validaties en interpretaties samenkomen. 
Geautomatiseerde tests, notebooks en duidelijke ADR’s (Architecture Decision Records) vormen daarbij een belangrijke basis om drift te beperken.

De belangrijkste bijdrage van deze bachelorproef is dus niet dat Python de oplossing op zichzelf is, maar dat het toont hoe een hybride, pragmatische en goed gedefinieerde aanpak de wetenschappelijke transparantie en technische betrouwbaarheid tegelijk kan verbeteren. 
De gekozen oplossing combineert Clean Architecture-principes, duidelijke domeinabstracties, Python-gebaseerde experimentatie en CI-geïntegreerde validatie. 
Hierdoor kunnen onderzoekers sneller experimenteren, sneller begrijpen wat er in het model gebeurt en beter formuleren welke assumpties werkelijk van invloed zijn op de uitkomsten.

Concluderend kan gesteld worden dat het onderzoek heeft aangetoond dat de “black box” van EMAV niet alleen technisch, maar vooral organisatorisch en communicatief wordt veroorzaakt. 
De technische complexiteit is realistisch, maar niet onoverkomelijk. 
De oplossing ligt in een goed afgestemde hybride architectuur waarin de vaste C\#-basis behouden blijft, terwijl Python en notebooks worden gebruikt als open, interactieve en reproduceerbare laag voor onderzoek en validatie. 
De belangrijkste stap voor ILVO is nu niet meer een complete herontwikkeling, maar het institutionaliseren van deze aanpak: docs-as-code, expliciete eigenaarschap, test- en notebook-integratie in de workflow en een bewuste cultuur van samenwerking tussen ontwikkelaars en onderzoekers.

Als deze aanbevelingen worden overgenomen, kan de organisatie profiteren van meer transparantie, minder documentation drift, beter reproduceerbare analyses en een kortere kennisoverdracht tussen technisch personeel en wetenschappelijk onderzoek. 
Daarmee wordt niet alleen de kwaliteit van de software verbeterd, maar ook de betrouwbaarheid en geloofwaardigheid van de onderzoeksresultaten zelf.

